\begin{proposition} \label{X.2.3}
Soit $X$ un pr\'esch\'ema localement noeth\'erien et soit $Y$ une partie ferm\'ee de $X$. Soit $\hat{X}$ le compl\'et\'e formel de $X$ le long de $Y$. Pour tout ouvert $U$ de $X$, $U\supset Y$, d\'esignons par $L_U$ (\resp $P_U$, \resp $E_U$) le foncteur qui \`a tout $\Oo_U$-Module coh\'erent localement libre (\resp \`a tout rev\^etement fini et plat de $U$, \resp \`a tout rev\^etement \'etale de $U$) associe son image inverse par $\hat{X} \to X$.
\begin{enumeratei}
\item
Si on~a $\Lef(X, Y)$, alors pour tout voisinage ouvert $U$ de $Y$, les foncteurs $L_U$, $P_U$ et $E_U$ sont pleinement fid\`eles.

\item
Si on~a $\Leff(X, Y)$, alors pour tout $\Oo_{\hat{X}}$-Module coh\'erent localement libre~$\Ee$ (\resp ...), il existe un ouvert $U$ et un $\Oo_U$-Module coh\'erent localement libre~$E$ (\resp ...), tels que $L_U(E) \simeq \Ee$ (\resp ...).
\end{enumeratei}
\end{proposition}
